\chapter{Время корреляции неравновесной ванны}
\label{ch:v10-nonequilibrium-bath-correlation-time}

Обратная максимальная частота
\begin{equation}
 \tau_{\rm UV}=\omega_{\rm UV}^{-1}
 =\frac{\ell_{\rm cell}}{2\sqrt3\,v_g}
\end{equation}
является естественной единицей времени. Но физическое время корреляции
определяется не только краем спектра, а всей нормированной функцией
корреляции:
\begin{equation}
 \tau_{\rm corr}=\int_0^\infty
 \frac{C(t)}{C(0)}\,dt.
\end{equation}

Положим $x=\omega_{\rm UV}t$. Два нормированных профиля с одной и той же
UV-шкалой дают
\begin{align}
 C_E(x)&=e^{-x}, & \tau_E\omega_{\rm UV}&=1,\\
 C_G(x)&=e^{-x^2}, & \tau_G\omega_{\rm UV}&=\frac{\sqrt\pi}{2}.
\end{align}
Они различаются также коротковременным наклоном:
$C_E'(0)=-1$, тогда как $C_G'(0)=0$.

Более того, непрерывное семейство
\begin{equation}
 C_a(x)=a e^{-x}+(1-a)e^{-x^2},\qquad 0\leq a\leq1,
\end{equation}
сохраняет $C_a(0)=1$ и ту же UV-единицу, но имеет
\begin{equation}
 \tau_a\omega_{\rm UV}
 =a+(1-a)\frac{\sqrt\pi}{2},
 \qquad C_a'(0)=-a.
\end{equation}
Следовательно, cutoff не выбирает ни интегральную память, ни её
коротковременный наклон.

Для заранее выбранного экспоненциального профиля родитель
\begin{equation}
 \mathcal P_\tau(r)=\frac12(r-1)^2,
 \qquad r=\tau_{\rm corr}\omega_{\rm UV},
\end{equation}
строго выбирает $r=1$. Но это условный родитель: сам выбор профиля в нём
уже содержится.

Размерная карта на
$(\tau_{\rm corr},\omega_{\rm UV},\ell_{\rm cell},v_g)$ имеет ранг/ядро
$2/2$. Якорь $v_g$ повышает ранг до $3$ и оставляет длино-временную орбиту;
только независимый якорь $\ell_{\rm cell}$ закрывает карту до $4/0$.

\begin{theorem}[Условное время экспоненциальной ванны]
После независимого выбора экспоненциального профиля корреляции UV-родитель
фиксирует $\tau_{\rm corr}=\omega_{\rm UV}^{-1}$ и тем самым
$\tau_{\rm corr}v_g/\ell_{\rm cell}=1/(2\sqrt3)$.
\end{theorem}

\begin{nogocriterion}[Cutoff не определяет память]
Экспоненциальный и гауссов профиль имеют одинаковую нормировку и одну
UV-шкалу, но разные интегральные времена и коротковременные наклоны.
Поэтому происхождение $\tau_{\rm corr}$ требует родителя полной спектральной
плотности, а не одной граничной частоты.
\end{nogocriterion}

\begin{protocol}\begin{enumerate}
\item условная архитектура времени корреляции: \textbf{$8/8$};
\item свидетель неединственности профиля: \textbf{$3/3$};
\item происхождение спектральной формы ванны: \textbf{$0/1$};
\item происхождение абсолютного времени корреляции: \textbf{$0/1$};
\item следующий гейт: \textbf{аудит спектральной плотности и масштаба памяти}.
\end{enumerate}\end{protocol}

Машинный сертификат сохранён в
\path{s2t_v10_cell_birth_four_volume_nonequilibrium_bath_correlation_time_parent_origin_gate_results.json}.